\section{From capacity to reconstruction}\label{sec:sufficiency}

Suppose $\theta$ is an admissible flow of finite mass $A_0>0$.
For a parameter $\eps>0$ to be chosen, set
\[
 \Theta=\eps\theta,\qquad M=\eps A_0,\qquad Q=1+A_0^2.
\]
Write $\Theta(v)$ for the incoming flow at $v\ne\rho$, put
$\Theta(\rho)=M$, and define $X_v=\Theta(v)/R_v$. Conservation gives
\begin{equation}\label{eq:profile-recursion}
 X_v=\sum_{w\in\ch(v)}r_eX_w.
\end{equation}

\subsection{A bounded profile and degraded observations}
Every positive-flow vertex extends to an infinite ray, by repeatedly
choosing a positive-flow child. Along a positive-flow edge,
$r_eX_w\leq X_v$, so
\[
 X_w^2-X_v^2\leq(1-r_e^2)X_w^2.
\]
Telescoping from the root to a vertex, and extending to a ray, yields
\begin{equation}\label{eq:profile}
 \sup_{v:\Theta(v)>0}X_v^2\leq M^2+\eps^2=\eps^2Q.
\end{equation}
Choose $\eps\sqrt Q\leq1$. Fix $n\geq1$. Independently pass
each observed spin at $v\in L_n$ through a channel with parameter
$\sqrt{X_v/3}$, using parameter zero when $X_v=0$.
Let $\widehat m_v,\widehat h_v$ be the information and absolute
third moment of the resulting descendant observation about
$\sigma_v$, for $|v|\leq n$. At the terminals,
\begin{equation}\label{eq:terminal}
 \widehat m_v=X_v,\qquad
 \widehat h_v=(X_v/3)^{3/2}.
\end{equation}
All these experiments belong to $\calC$. Subadditivity from
\eqref{eq:saving} and \eqref{eq:profile-recursion} gives inductively
$\widehat m_v\leq X_v$. Zero-flow branches are uninformative and
may be omitted. All computations at depth $n$ involve finitely many
vertices, because $T$ is locally finite.

At a positive-flow vertex define
\begin{equation}\label{eq:splits}
 t_{vw}=\frac{\Theta(w)}{\Theta(v)},\qquad
 \beta_v=1-\sum_w t_{vw}^2,\qquad
 U_v=\frac{\widehat h_v}{X_v}.
\end{equation}
The sums are over its positive-flow children, and $\sum_w t_{vw}=1$.

\begin{lemma}\label{lem:deficit}
For every positive-flow vertex $v$ with $|v|<n$, put
\[
 D_v=\sum_w r_e\widehat m_w-\widehat m_v.
\]
With the universal constant $C_0=2000$,
\begin{align}
 U_v&\leq\sum_wt_{vw}\lambda_eU_w+C_0X_v\beta_v,
                       \label{eq:H}\\
 0\leq\frac{D_v}{X_v}
 &\leq C_0\left[X_v\sum_wt_{vw}(1-t_{vw})\lambda_eU_w
                     +X_v^2\beta_v\right].\label{eq:D}
\end{align}
\end{lemma}

\begin{proof}
Order the children arbitrarily and write
$b_i=r_e\widehat m_w$, $h'_i=\lambda_e^3\widehat h_w$ for
their channel-transformed moments. Set $S=\sum_i b_i$ and
$P_2=\sum_{i<j}b_ib_j$. Then
\begin{equation}\label{eq:pair-bound}
 b_i\leq t_iX_v,\qquad S\leq X_v,\qquad
 P_2\leq\frac{X_v^2\beta_v}{2}.
\end{equation}
Sequential application of \eqref{eq:local-h}, using subadditivity
for each partial product, gives
\[
 \widehat h_v\leq\sum_i h'_i+460P_2.
\]
Since $h'_i=t_iX_v\lambda_eU_w$, this proves \eqref{eq:H}.

By \eqref{eq:local-m}, the deficit at each multiplication is at
most seven times the sum of the two cross products of a second and
an absolute third moment. Summing over the ordered products yields
\[
 D_v\leq7\sum_i h'_i(S-b_i)
              +7\cdot460\sum_{i<j<k}b_ib_jb_k.
\]
The triple sum is at most $X_vP_2$. Also
$S-b_i=\sum_{j\ne i}b_j\leq X_v(1-t_i)$.
After division by $X_v$, \eqref{eq:pair-bound} bounds the
triple term by $1610X_v^2\beta_v$, while the pair term is
at most $7X_v\sum_i t_i(1-t_i)\lambda_eU_w$.
Subadditivity gives $D_v\geq0$. Both coefficients are bounded
by $C_0$. With one positive-flow child, $\beta_v=D_v=0$.
\end{proof}

\subsection{Energy from branch splitting}
Choose a terminal $v\in L_n$ with probability $\Theta(v)/M$,
and let $V_0=\rho,V_1,\ldots,V_n=v$ be its ancestral path.
Denote expectation for this choice by $\mathbf E$; it is separate
from the expectations defining posterior moments. Let
$\mathcal F_i=\sigma(V_0,\ldots,V_i)$ and write
\[
 X_i=X_{V_i},\quad t_{i+1}=t_{V_iV_{i+1}},\quad
 \alpha_i=1-t_{i+1},
\]
with $\lambda_{i+1},r_{i+1}$ the parameters of that edge.
Conditional on $\mathcal F_i$, the next child is chosen with
probabilities $t_{V_iw}$, so
$\mathbf E[\alpha_i\mid\mathcal F_i]=\beta_{V_i}$.

\begin{lemma}\label{lem:energy}
For every $n$,
\begin{equation}\label{eq:energy}
 \mathbf E\sum_{i=0}^{n-1}\alpha_iX_i^2\leq\eps^2Q.
\end{equation}
\end{lemma}
\begin{proof}
Using $r_eX_w=t_{vw}X_v$, direct conditional expectation gives
\begin{align*}
 X_i^2-\mathbf E[X_{i+1}^2\mid\mathcal F_i]
 +\mathbf E[(1-r_{i+1}^2)X_{i+1}^2\mid\mathcal F_i]
 &=X_i^2\left(1-\sum_wt_{V_iw}^3\right)\\
 &\geq X_i^2\beta_{V_i}.
\end{align*}
Sum over $i$ and take expectations. The first two terms telescope
to $M^2-\mathbf E X_n^2$. Every sampled path extends to an infinite
ray, so its sum of $(1-r_{i+1}^2)X_{i+1}^2$ is at most
$\eps^2$ by admissibility. Dropping the negative terminal term
and replacing $\beta_{V_i}$ by its conditional expectation proves
\eqref{eq:energy}.
\end{proof}

For $0\leq i<j\leq n$ put
\[
 K_{ij}=\prod_{k=i+1}^j\sqrt{t_k},\qquad K_{ii}=1.
\]
On every sampled path,
\begin{equation}\label{eq:kernel}
 \sum_{j=i+1}^{n-1}\alpha_jK_{ij}\leq2,
 \qquad
 \sum_{i=0}^{j-1}\alpha_iK_{ij}\leq2\quad(1\leq j\leq n).
\end{equation}
Indeed, $1-t\leq2(1-\sqrt t)$ for $0\leq t\leq1$.
For the first bound use
$\alpha_jK_{ij}\leq2(K_{ij}-K_{i,j+1})$ and telescope.
For the second use
$\alpha_iK_{ij}\leq2(K_{i+1,j}-K_{ij})$ and telescope.
In particular, the second bound includes the terminal index $j=n$.

\begin{lemma}\label{lem:double}
On every sampled path,
\begin{equation}\label{eq:double}
 \sum_{0\leq i<j<n}\alpha_i\alpha_jX_iX_j
                       \prod_{k=i+1}^j\lambda_k
 \leq2\sum_{i=0}^{n-1}\alpha_iX_i^2.
\end{equation}
\end{lemma}
\begin{proof}
The profile relation gives
$X_j=X_i\prod_{k=i+1}^j(t_k/\lambda_k^2)$. Hence
\[
 X_iX_j\prod_{k=i+1}^j\lambda_k
 =X_i^{3/2}X_j^{1/2}K_{ij}
 \leq\left(\frac34X_i^2+\frac14X_j^2\right)K_{ij}.
\]
Multiply by $\alpha_i\alpha_j$ and sum. The row and column
bounds in \eqref{eq:kernel} bound the two terms by, respectively,
$\tfrac32\sum_i\alpha_iX_i^2$ and
$\tfrac12\sum_j\alpha_jX_j^2$.
\end{proof}

\subsection{The total information loss}
Unfold \eqref{eq:H} down to level $n$ and insert it into
\eqref{eq:D}. Taking path expectations gives
\begin{align}
 \mathbf E\sum_{i<n}\frac{D_{V_i}}{X_i}
 \leq{}& C_0\eps^2Q\notag\\
 &+C_0^2\mathbf E\sum_{i<j<n}\alpha_i\alpha_jX_iX_j
                          \prod_{k=i+1}^j\lambda_k\notag\\
 &+C_0\mathbf E\sum_{i<n}\alpha_iX_i
                \left(\prod_{k=i+1}^n\lambda_k\right)U_{V_n}.
                \label{eq:unfolded}
\end{align}
To justify the conditional-expectation step explicitly,
\eqref{eq:D} first contributes
$X_i\mathbf E[\alpha_i\lambda_{i+1}U_{V_{i+1}}\mid\mathcal F_i]$.
Iteration of \eqref{eq:H} from $i+1$ replaces $U_{V_{i+1}}$
by its terminal contribution and the sum of
$C_0X_j\beta_{V_j}$ with the intervening channel products.
For $i<j$, its multiplier
$\alpha_iX_iX_j\prod_{k=i+1}^j\lambda_k$ is
$\mathcal F_j$-measurable, so $\beta_{V_j}$ may be replaced by
$\alpha_j$ inside expectation. The direct $X_i^2\beta_{V_i}$
term is bounded by Lemma~\ref{lem:energy}. This proves
\eqref{eq:unfolded} using only finite sums.

By Lemmas~\ref{lem:energy} and~\ref{lem:double}, the middle
term in \eqref{eq:unfolded} is at most $2C_0^2\eps^2Q$.
At the terminal, \eqref{eq:terminal} gives
$U_{V_n}=\sqrt{X_n}/(3\sqrt3)$, and therefore
\[
 X_i\left(\prod_{k=i+1}^n\lambda_k\right)U_{V_n}
       =\frac{X_i^{3/2}}{3\sqrt3}K_{in}.
\]
Use \eqref{eq:profile} and the terminal case of
\eqref{eq:kernel} to bound its sum with $\alpha_i$ by
$2\eps^{3/2}Q^{3/4}/(3\sqrt3)$. It follows that
\begin{equation}\label{eq:loss}
 \mathbf E\sum_{i<n}\frac{D_{V_i}}{X_i}
 \leq(C_0+2C_0^2)\eps^2Q
             +\frac{2C_0}{3\sqrt3}\eps^{3/2}Q^{3/4}.
\end{equation}
Choose $\eps=10^{-4}/\sqrt Q$. Then the right side is
\[
 (2000+2\cdot2000^2)10^{-8}
       +\frac{4000}{3\sqrt3}10^{-6}
 =0.08002+\frac1{750\sqrt3}<\frac12.
\]
This choice also satisfies $\eps\sqrt Q\leq1$.

Finally, multiply the definition of $D_v$ by $R_v$ and sum
over the finite truncation at depth $n$. The intermediate terms
cancel, and \eqref{eq:terminal} gives
\begin{align*}
 \widehat m_\rho
 &=\sum_{|v|=n}R_vX_v-\sum_{|v|<n}R_vD_v\\
 &=M-M\mathbf E\sum_{i<n}\frac{D_{V_i}}{X_i}
 \geq\frac M2.
\end{align*}
Zero-flow vertices contribute zero, so the displayed ratios are only
taken on sampled paths. Every level has total flow $M$.
The original level observation is more informative than the degraded
one. Therefore, for every $n\geq1$,
\[
 m_n\geq\frac{\eps A_0}{2}
       =\frac1{20000}\frac{A_0}{\sqrt{1+A_0^2}}.
\]
Pass to the limit in $n$, then take admissible flow masses tending
to $C(T)$, or to infinity when $C(T)=\infty$. This proves the lower
bound in \eqref{eq:main-bounds}. If $C(T)=0$, that lower bound is
automatic. Together with \eqref{eq:strong-upper} and
\eqref{eq:tv}, the proof of Theorem~\ref{thm:main} is complete.
